% Generated by scripts/analysis/revision_configuration_tables.py; no sampling.
\section{紧凑研究的最终目标与配置}
\label{sec:compact-final-config}
本节直接读取正式紧凑协议及其冻结目标目录，未使用历史CPU配置替代本轮定义。全部核和两设备共用每目标的固定仿射坐标。记基础模型坐标为$r=c+Lq$，其中$q$是实际采样坐标；原参数为$\theta=T_0(r)$。除H2外$T_0$为恒等映射，H2再将非中心化变量转换为漏斗原尺度。变换后的密度含$\log|\det L|$。

\begin{table}[htbp]\centering\small
\begin{tabular}{lrrrl}\toprule
目标 & 维数 & MALA $h$ & RWM $s$ & 固定仿射坐标\\\midrule
G1 & 8 & 0.8 & 0.848528137 & $c=0,L=I$\\
G2 & 64 & 0.4 & 0.3 & $c=0,L=\operatorname{chol}(\Sigma)$\\
A1 & 64 & 0.4 & 0.3 & $c=m,L=\operatorname{chol}(\Sigma)$\\
L1 & 8 & 0.8 & 0.848528137 & 众数及观测信息\\
L2 & 8 & 0.8 & 0.848528137 & 众数及观测信息\\
H1 & 9 & 0.480749857 & 1.66666667 & $c=0,L=I$\\
H2 & 9 & 0.769199771 & 0.8 & $c=0,L=I$\\
M1 & 8 & 0.8 & 0.848528137 & $c=0,L=I$\\
W1 & 2 & 1.26992084 & 1.69705627 & 众数及观测信息\\
\bottomrule\end{tabular}
\caption{正式紧凑研究采用的最终配置。$h$遵循正文噪声$\sqrt{2h}$的MALA约定；$s$是RWM提议标准差。表内数字为显示精度，完整浮点数、中心和矩阵保存在随附配置JSON及原协议中。}
\label{tab:compact-final-steps}\end{table}

\paragraph{高斯与AR(1)目标。}
G1为$N_8(0,I)$。G2为$N_{64}(0,U\Lambda U^\top)$，$\Lambda_{jj}=100^{(j-1)/63}$；$U$由固定$64\times64$标准正态矩阵的QR分解取得。原始协方差条件数100，但完整Cholesky变换使采样坐标为$N_{64}(0,I)$。A1的生成模型为$x_1\sim N(0,1)$、$x_t=0.9x_{t-1}+\sqrt{1-0.9^2}\epsilon_t$、$y_t=x_t+\eta_t$，$t=1,\ldots,64$，噪声相互独立且为标准正态。记$K_{ij}=0.9^{|i-j|}$，固定数据后的后验为$N(m,\Sigma)$，$\Sigma=(K^{-1}+I)^{-1}$、$m=\Sigma y$。A1也使用完整Cholesky变换，因而采样坐标同为标准正态。G2/A1的原尺度相关结构不能直接作为本轮核差异的原因；维数、各核移动效率、稠密目标计算及有限精度仍影响被测工作流。

\paragraph{回归目标。}
L1/L2均为8参数logistic回归，样本数分别为256/4096，先验$\beta_j\overset{\mathrm{iid}}\sim N(0,2.5^2)$。设计矩阵第一列为1，其余七列从标准正态生成后按每列样本均值和总体式标准差（分母$n$）标准化。生成系数来自$N_8(0,2.5^2I)$且两模型共用，观测独立来自$\operatorname{Bernoulli}(\sigma(X_i\beta))$。W1采用3020条真实水井记录，$\Pr(y_i=1)=\sigma(\alpha+\beta d_i/100)$，先验在$(\alpha,\beta)$上平坦。L1/L2/W1以零点开始的trust-exact优化求众数$c$，接受条件为目标有限及解析梯度最大绝对值不超过$10^{-7}$；$L=\operatorname{chol}(H(c)^{-1})$。L1/L2的$H=X^\top\operatorname{diag}\{p_i(1-p_i)\}X+I/2.5^2$，W1不加先验项。优化最多300步，未做特征值修复，未使用参考后验样本；这只是固定坐标选择，不将Laplace近似当作参考真值。

\paragraph{漏斗与双峰目标。}
H1为$v\sim N(0,3^2)$，$x_i\mid v\sim N(0,e^v)$，$i=1,\ldots,8$。H2采样$v\sim N(0,3^2),z_i\sim N(0,1)$并输出$x_i=z_i e^{v/2}$。二者基础坐标不同，原参数后验对应。M1为$\tfrac12N_8(-5e_1,I)+\tfrac12N_8(5e_1,I)$，即$\delta=5$，未额外平移或旋转，因此本例$q=\theta$数值相同，正文仍统一用$\theta$表示原参数。

\paragraph{固定数据的生成与定位。}
合成目标使用NumPy默认生成器和\texttt{SeedSequence}。G2、A1和logistic生成系数分别从地址$[20261003,11]$开始；G2先生成旋转矩阵，A1先生成潜状态创新再生成观测噪声。L1/L2的设计及二元观测分别使用地址$[20261003,12,n]$。相同地址在不同模型内重新开始，并不表示G2与A1共用同一目标数组。归档的模型数据和协方差是执行依据，完整再生成核对最大绝对差小于$10^{-11}$。无需依赖本机目录即可在正式证据包的\texttt{source/benchmark/protocols/windows-native-v1.json}及\texttt{catalog.json}定位完整定义；配置JSON还给出来源SHA及完整$c,L$。

\paragraph{开发选择、初值与核控制。}
坐标构造规则预先固定，随后在独立开发输入上分别选择MH步长，而非从正式结果重新调参。RWM候选为$\{0.5,1,1.5,2.4,3.5,5\}/\sqrt d$，MALA候选为$\{0.1,0.3,0.6,1,1.6,2.4\}/d^{1/3}$；每候选两份四链输入，丢弃512步后保留1024步，以两次均有效时的平均ESJD/d最大者入选，平局取较小步长。正式四链共用原件中的$q_0\sim N(0,4I)$，M1首坐标替换为$(-5,5,-5,5)$。MH和NUTS读取同一初值，原尺度初值为$T_0(c+Lq_0)$；NUTS另用独立随机流和适应预热，并非与MH共享路径。216份实际初值逐数组核对通过。

正式保留预算为1024/4096步；MH先丢弃512步，窗口32，quasi-DEER每窗口最多2048轮，Picard最多$N+512$轮。CPU Pyro NUTS采用四个spawn子进程，每进程一链、一条torch计算线程和一条interop线程；每链预热1024步、目标接受率0.8、最大树深8，适应步长及对角质量矩阵（\texttt{full\_mass=False}），未使用JIT。各链适应后的参数保存在输出中，而非用表内MH步长替代。

\paragraph{资源与失败规则。}
主MH工作进程设置torch计算线程4、interop线程1，CUDA任务也保留这部分主机执行。3888份主进程环境记录均为上述设置，浮点数为float64；OMP/MKL/OpenBLAS/VECLIB线程环境变量未设置，因此不将torch设置解释为所有库的统一线程硬上限。线程数也不等于独占物理核心预算。每任务存储/工作区估计上限2048 MiB，进程树RSS监督上限8 GiB，原生Windows Job提交内存上限12 GiB；它们是不同计量规则，2048 MiB不是通用分配器硬上限。任务启动保留可用主存至少12 GiB及磁盘至少4 GiB，CUDA工作进程另检查GPU空闲至少3 GiB，运行磁盘底线1 GiB。边界资源不足则暂停启动，不把未运行任务标为数值失败；进程异常、非有限值、迭代耗尽和路径标准失配均保留。没有总实验时限、数值失败重试或顺序回退来改变样本资格。12 GiB限制本身不证明所见进程失败均由内存耗尽导致。

\paragraph{版本依据。}
以上读取的是\texttt{windows-compact-inference-v1}，冻结执行源\texttt{3a37a89}，协议摘要前缀\texttt{91bbff0ec47f}。冻结文件和依赖版本保持不变；本次仅导出配置、核对既有输入，新增采样和调参均为零。
